\subsection{General Anomaly Detection Methods and Foundations}

In the field of anomaly detection, the paper "Anomaly Detection: A Survey" \cite{paper_AnomalyDetection_Northwestern} serves as a central methodological foundation. It provides a systematic classification of approaches, including cluster-based and statistical methods, and explains their respective strengths, weaknesses, and typical use cases. Furthermore, it differentiates between various types of anomalies, such as point anomalies and contextual anomalies, and discusses the most suitable detection strategies for each category. This makes it an indispensable reference for modern anomaly detection procedures.\\
Building on this established methodological framework, a solid basis is created upon which more recent approaches, particularly deep learning models and methods tailored specifically to EVSE, can be developed. This also includes the LSTM-based anomaly detection presented in this work. Therefore, the paper highlights both the breadth and depth of anomaly detection research and ensures that specialized techniques are embedded within a well-founded conceptual context \cite{paper_AnomalyDetection_Northwestern}.\\
\\
Another relevant contribution in the broader context of cybersecurity is presented by Rahman et al. (2024), titled "ChronoCTI: Temporal Pattern Extraction from Cyber Threat Intelligence Reports" \cite{attack_over_time}. The authors propose an automated pipeline that leverages modern machine learning techniques and natural language processing to extract recurring, time-based sequences of malicious activities from unstructured reports. In contrast to traditional approaches that primarily focus on static indicators such as IP addresses or malware signatures, this work emphasizes the analysis of how an attack evolves over time.\\
The findings demonstrate that incorporating the temporal progression of malicious activity into detection processes can significantly enhance threat identification. For EVSE cybersecurity, this perspective is particularly valuable, as certain attacks on charging infrastructure may unfold gradually, with patterns becoming apparent only when their temporal dimension is considered. Such an approach aligns closely with the strengths of LSTM-based anomaly detection, which is well-suited to capturing dependencies over time and can thereby improve the detection of subtle, long-running attacks.
